% Image credits for the photographs and AI illustrations of Book 3 (University, Year 1), French edition.
% Machine-readable license records live in images/book3/CREDITS.md; keep this
% file in sync with frontmatter/image-credits-book3.tex (one \item per item there).
% Every credit, author, dataset and licence is kept exactly as in the English file.
\chapter*{Crédits des images}

Les dessins, cartes et graphiques de ce livre sont originaux~; les cartes
et graphiques tracés à partir de données publiées en citent la source dans
la légende ou dans cette liste. Les photographies sont utilisées sous leurs
licences libres respectives, avec nos remerciements à leurs auteurs~; les
liens complets vers les sources et les adresses des licences de chaque
photographie figurent dans \texttt{images/book3/CREDITS.md}, dans le dépôt
du livre.

\bigskip
À côté des photographies et des dessins originaux, certaines figures sont
des illustrations générées par IA, produites avec le générateur d'images
d'OpenAI à partir de consignes rédigées par les éditeurs du livre, et
vérifiées pour leur exactitude scientifique avant d'être retenues. Elles
ne montrent que des scènes génériques, jamais un lieu réel, une personne
réelle ou un spécimen réel. La consigne complète de chaque image générée
est consignée dans \texttt{images/book3/ai/PROMPTS.md}, dans le dépôt.

\bigskip
\noindent Photographies et données, par chapitre~:
\begin{itemize}\small
  \item Chap.~1 (astronomie sphérique)~: traînées d'étoiles au-dessus de l'observatoire du Paranal, R.~Wesson/ESO, CC BY 4.0.
  \item Chap.~2 (télescopes, détecteurs et photométrie)~: six segments du miroir primaire du télescope spatial James Webb avant les essais cryogéniques, NASA/MSFC/David Higginbotham, domaine public.
  \item Chap.~3 (planétologie comparée)~: les planètes telluriques à l'échelle, NASA~; Vénus au radar de Magellan, NASA/JPL~; Mars vue par l'orbiteur Viking~1, NASA/USGS~; Europe, NASA/JPL/DLR~; Io, NASA/JPL/University of Arizona~; tous domaine public.
  \item Chap.~4 (les cartes géologiques)~: William Smith, \emph{A Delineation of the Strata of England and Wales, with part of Scotland}, 1815~; domaine public (Wikimedia Commons).
  \item Chap.~5 (les coupes géologiques)~: la discordance de Hutton à Siccar Point, PawełS, CC BY-SA 4.0 (Wikimedia Commons).
  \item Chap.~6 (méthodes de terrain et de laboratoire~I)~: un granite à biotite en lame mince entre polariseurs croisés, Darla Sondrol, CC0 (Wikimedia Commons).
  \item Chap.~7 (les minéraux des roches)~: échantillons d'olivine, de quartz, de biotite, d'almandin et de hornblende, Robert M.~Lavinsky, CC BY-SA 3.0~; labradorite, Twyla Baker, CC0~; rhomboèdres de clivage de calcite, Giovanna Canu et Eva Santini, CC BY 4.0 (tous Wikimedia Commons).
  \item Chap.~8 (minéralogie optique)~: un gabbro en lame mince en lumière polarisée non analysée et entre polariseurs croisés, Chmee2, CC BY-SA 3.0~; Auguste Michel-Lévy étudiant, domaine public (Wikimedia Commons). Les teintes de polarisation sont calculées à partir des fonctions colorimétriques CIE 1931 et de l'illuminant standard CIE D65, \copyright{} CIE, CC BY-SA 4.0.
  \item Chap.~9 (la Terre en couches)~: Andrija Mohorovičić, domaine public~; Inge Lehmann en 1932, photographie d'Even Neuhaus, domaine public (Wikimedia Commons). PREM~: A.~M.~Dziewonski et D.~L.~Anderson (1981), tel que tabulé dans ObsPy.
  \item Chap.~10 (le bilan thermique de la Terre)~: flux géothermique mondial d'après J.~H.~Davies et D.~R.~Davies (2010), \emph{Solid Earth}, CC BY 3.0~; adiabate du manteau d'après T.~Katsura (2022), \emph{JGR Solid Earth}, CC BY.
  \item Chap.~11 (isostasie et anomalies de pesanteur)~: l'appareil à pendules de Vening Meinesz, Broodt, CC BY-SA 4.0 (Wikimedia Commons). Altitudes d'après ETOPO1 (NOAA National Centers for Environmental Information), échantillonnées à un degré.
  \item Chap.~12 (cinématique des plaques sur une sphère)~: une station permanente de positionnement par satellite, Alpsdake, CC BY-SA 4.0 (Wikimedia Commons). Limites de plaques d'après PB2002 (P.~Bird, 2003~; GeoJSON de H.~Ahlenius, ODC-By 1.0)~; côtes d'après Natural Earth (domaine public).
  \item Chap.~13 (rifts, dorsales et marges passives)~: fissures ouvertes à Þingvellir, Islande, Olga Ernst, CC BY-SA 4.0~; laves en coussins sur le rift des Galápagos, NOAA Okeanos Explorer Program, CC BY 2.0 (tous deux Wikimedia Commons). Profondeurs d'après ETOPO1 (NOAA National Centers for Environmental Information) et âges du plancher océanique d'après M.~Seton et ses collègues (2020), EarthByte, CC BY 4.0.
  \item Chap.~14 (les marges actives)~: éruption du pic Sarychev vue depuis la Station spatiale internationale, NASA, domaine public (Wikimedia Commons). Séismes d'après le USGS ANSS Comprehensive Catalog (domaine public)~; volcans d'après le Global Volcanism Program de la Smithsonian Institution~; limites de plaques d'après PB2002 (P.~Bird, 2003~; GeoJSON de H.~Ahlenius, ODC-By 1.0)~; côtes d'après Natural Earth (domaine public)~; âges du plancher océanique d'après M.~Seton et ses collègues (2020), EarthByte, CC BY 4.0.
  \item Chap.~15 (points chauds et volcanisme intraplaque)~: une fontaine de lave à Pu\textquoteleft u \textquoteleft O\textquoteleft o, Kilauea, G.~E.~Ulrich, U.S. Geological Survey, domaine public~; les Ghâts occidentaux à Matheran, Nicholas (Nichalp), CC BY-SA 2.5 (tous deux Wikimedia Commons). Positions des points chauds, azimuts et vitesses des traces d'après le tableau~1 du supplément électronique de W.~J.~Morgan et J.~P.~Morgan (2007)~; profondeurs d'après ETOPO1 (NOAA National Centers for Environmental Information)~; limites de plaques d'après PB2002 (P.~Bird, 2003~; GeoJSON de H.~Ahlenius, ODC-By 1.0)~; côtes d'après Natural Earth (domaine public).
  \item Chap.~16 (contrainte, déformation et rhéologie)~: une cellule triaxiale à haute pression pour mesurer la résistance des roches, Darryl Peroni, CSIRO, CC BY 3.0 (Wikimedia Commons). Paramètres de fluage par dislocations du quartzite humide (Gleason et Tullis, 1995) et de l'olivine sèche (Hirth et Kohlstedt, 2003) tels que compilés par M.~A.~Lopez-Sanchez dans \emph{strength\_envelopes} (MPL-2.0).
  \item Chap.~17 (failles, plis et fabriques)~: un miroir de faille à Corona Heights, San Francisco, Dextersinister, CC BY-SA 4.0~; brèche de faille du chevauchement de Keystone, Nevada, Jstuby, domaine public~; plis en chevrons à Millook, Cornouailles, Smalljim, CC BY 3.0~; un anticlinal à clivage de plan axial près de Potes, Cantabrie, Mikenorton, CC BY-SA 3.0 (tous Wikimedia Commons).
  \item Chap.~18 (la source sismique)~: terrain et clôtures séparés par la rupture de 1906 dans le comté de Marin, \emph{Popular Science Monthly} vol.~69 (1906), domaine public~; Wallace Creek décalé à travers la faille de San Andreas, R.~E.~Wallace, U.S. Geological Survey, domaine public (tous deux Wikimedia Commons). Séismes d'après le USGS ANSS Comprehensive Catalog (domaine public)~; le mécanisme au foyer du séisme de Tohoku de 2011 d'après le tenseur de moment en phase W du USGS (domaine public).
  \item Chap.~19 (l'altération et les argiles)~: granite pourri et sable granitique à Targassonne, Pyrénées, Pinpin, CC~BY~2.5~; bauxite sur grès kaolinique à Weipa, Australie, Werner Schellmann, CC~BY-SA~2.5~; cristaux de kaolinite au microscope électronique à balayage, U.S. Geological Survey, domaine public (tous Wikimedia Commons). Rayons ioniques d'après R.~D.~Shannon (1976).
  \item Chap.~20 (érosion et transport)~: le panache de crue du Mississippi vu par l'instrument MODIS de Terra, 4~mars 2018, image de Joshua Stevens, NASA Earth Observatory, domaine public (Wikimedia Commons).
  \item Chap.~21 (les sédiments clastiques et leurs milieux)~: stratification entrecroisée dans le grès d'Aztec, Nevada, Chip Carroon, Bureau of Land Management, CC~BY-SA~2.0~; le delta de la Léna vu par Landsat~7, U.S. Geological Survey, domaine public~; une turbidite montrant les divisions de la séquence de Bouma, Becke-Oese, Jo Weber, CC~BY-SA~3.0 (tous Wikimedia Commons).
  \item Chap.~22 (les sédiments carbonatés, siliceux et évaporitiques)~: Bora Bora vue depuis la Station spatiale internationale, NASA Johnson Space Center, domaine public~; une cellule d'\emph{Emiliania huxleyi}, Alison R.~Taylor, University of North Carolina Wilmington Microscopy Facility, CC~BY~2.5~; \emph{Orbulina universa}, H.~Takagi, K.~Kimoto, T.~Fujiki, H.~Saito, C.~Schmidt, M.~Kucera et K.~Moriya, CC~BY~4.0~; diatomées du détroit de McMurdo, Gordon T.~Taylor, Stony Brook University, NOAA Corps Collection, domaine public (toutes Wikimedia Commons).
  \item Chap.~23 (la diagenèse et les bassins sédimentaires)~: surcroissances de quartz dans le grès de Potsdam, M.~C.~Rygel, CC~BY-SA~4.0 (Wikimedia Commons). Empilement du niveau de la mer de R.~M.~Spratt et L.~E.~Lisiecki (2016), NOAA NCEI World Data Service for Paleoclimatology~; paramètres de compaction de J.~G.~Sclater et P.~A.~F.~Christie (1980) tels que tabulés dans pyBacktrack (EarthByte).
  \item Chap.~24 (l'échelle stratigraphique)~: clou d'or du GSSP du Sélandien à Zumaia, Josugoni, CC~BY-SA~4.0 (Wikimedia Commons)~; la limite Crétacé--Paléogène au parc d'État du lac Trinidad, Nationalparks, CC~BY-SA~2.5 (Wikimedia Commons). Âges des bases et couleurs de la charte chronostratigraphique internationale, International Commission on Stratigraphy, CC~BY~4.0.
  \item Chap.~25 (les méthodes de datation radiométrique)~: un grain de zircon en cathodoluminescence, Emmanuel Roquette, CC~BY-SA~4.0~; un spectromètre de masse à ionisation thermique, Radiogenic, CC0 (tous deux Wikimedia Commons). La courbe d'étalonnage IntCal20 de P.~J.~Reimer et ses collègues (2020), \emph{Radiocarbon} 62, CC~BY~4.0, telle que distribuée avec le paquet rintcal (M.~Blaauw).
  \item Chap.~26 (la structure de l'atmosphère et de l'océan)~: remontée d'un CTD et de sa rosette dans le détroit de Béring, NASA Goddard Space Flight Center, domaine public (Wikimedia Commons). Température et salinité du World Ocean Atlas 2023 (NOAA National Centers for Environmental Information).
  \item Chap.~27 (le rayonnement et le transport de chaleur)~: la Terre éclairée vue par le Deep Space Climate Observatory, NASA Goddard Space Flight Center, domaine public (Wikimedia Commons). Rayonnement, flux de surface et précipitations de la Réanalyse~1 NCEP-NCAR (E.~Kalnay et ses collègues, 1996), données fournies par le NOAA PSL, Boulder, Colorado, États-Unis, depuis leur site https://psl.noaa.gov.
  \item Chap.~28 (la circulation générale de l'atmosphère)~: un cyclone extratropical sur le golfe d'Alaska, NASA MODIS/LANCE, domaine public~; la zone de convergence intertropicale d'après des données de nuages de GOES-11, NASA Earth Observatory, domaine public (tous deux Wikimedia Commons). Vents et pression au niveau de la mer de la Réanalyse~1 NCEP-NCAR (E.~Kalnay et ses collègues, 1996), données fournies par le NOAA PSL, Boulder, Colorado, États-Unis, depuis leur site https://psl.noaa.gov.
  \item Chap.~29 (la circulation de l'océan)~: le Gulf Stream vu à travers la température de la surface de la mer, NOAA NESDIS, domaine public~; un flotteur profileur profond en cours de mise à l'eau, photographe inconnu de la NOAA, domaine public (tous deux Wikimedia Commons).
\end{itemize}
